\documentclass[11pt]{article}
\usepackage[margin=1in]{geometry}
\usepackage{times}
\usepackage{latexsym}
\usepackage{amsmath}
\usepackage{amssymb}
\usepackage{graphicx}
\usepackage{booktabs}
\usepackage{multirow}
\usepackage{xcolor}
\usepackage{url}
\usepackage{hyperref}
\usepackage{amssymb}
\usepackage{pifont}
\renewcommand{\checkmark}{\ding{51}}

\title{FREIGE: Faithful and Evidence-Grounded Document-Level Relation Extraction with LLMs}

\author{Anonymous}

\begin{document}
\maketitle

\begin{abstract}
Large language models achieve near-encoder-SOTA extraction accuracy on document-level relation extraction (DocRE), yet no existing method produces verifiable evidence for extracted relations. We formalize \textbf{evidence-grounded DocRE}---simultaneously generating relation triples and supporting evidence sentences---and systematically compare online RL (GRPO), offline preference optimization (DPO), and rejection sampling fine-tuning (RSFT) for this task.
Our key finding is a \textbf{decoupling spectrum}: RL optimization pressure correlates positively with evidence-extraction decoupling severity. GRPO with QLoRA induces catastrophic decoupling (rel\_F1$<$0.02), GRPO with bf16 shows mild decoupling ($-$2pp), DPO exhibits gradual decoupling across epochs, while RSFT avoids decoupling entirely. This pattern holds across model families (Qwen3, LLaMA-3.1) and scales (1.7B--8B).
We propose Contrastive Evidence Discrimination (CED) as an NLI-based scoring function for RSFT and introduce the Evidence-Decoupling Calibration Ratio (EDCR) as a diagnostic metric validated against human judgments ($\rho$=0.78). On the DocRED test set, our RSFT approach improves Ign-F1 by +4.8pp over SFT while maintaining evidence quality, and the CED reranker achieves 46.5 Ign-F1 with Evi-F1=81.2---recovering 97\% of the evidence generation cost with faithful evidence.
\end{abstract}

\section{Introduction}

Document-level relation extraction (DocRE) requires identifying relations between entities across multiple sentences within a document. While encoder-based approaches~\cite{tan2022dreeam,xu2022sais} achieve strong performance, recent work has shown that generative LLMs can approach encoder SOTA on extraction accuracy. However, a critical gap remains: no generative LLM method produces verifiable evidence---the supporting sentences that justify each extracted relation.

We argue that evidence-grounded predictions are essential for deployment, where users need to verify model outputs. Encoder-based methods like SAIS and DREEAM perform evidence \emph{retrieval}---classifying which existing sentences support a relation---which serves as auxiliary supervision to improve RE. In contrast, we study generative evidence \emph{production}: requiring the LLM to generate evidence text alongside relation triples, which imposes an additional generation burden.

This burden is substantial. Our experiments reveal an \textbf{evidence generation tax} of approximately 7pp Ign-F1: an SFT model without evidence generation achieves 47.2 Ign-F1, while the same model trained to produce evidence drops to 39.7. The central question becomes: \emph{can we recover this cost while maintaining faithful evidence?}

We systematically explore three optimization paradigms:
\begin{itemize}
\item \textbf{Online RL (GRPO)}: Group Relative Policy Optimization with CED-based reward
\item \textbf{Offline Preference (DPO)}: Direct Preference Optimization on CED-ranked pairs
\item \textbf{Offline RL (RSFT)}: Rejection Sampling Fine-Tuning with NLI-based quality filtering
\end{itemize}

Our investigation reveals a striking \textbf{decoupling spectrum}: as optimization pressure increases, evidence quality improves but extraction accuracy degrades---a phenomenon we term \emph{evidence-extraction decoupling}. QLoRA quantization amplifies this to catastrophic levels (rel\_F1$<$0.02), establishing that decoupling is an intra-task multi-objective trade-off intrinsic to RL-for-IE, not a model-specific artifact.

RSFT emerges as the only approach that improves both extraction and evidence quality simultaneously. We validate our findings across three dimensions: cross-scale (Qwen3-1.7B/4B/8B), cross-dataset (DocRED, Re-DocRED), and cross-family (Qwen3, LLaMA-3.1), achieving consistent improvements of +3.3--5.1pp Ign-F1 over SFT baselines.

Our contributions are:
\begin{enumerate}
\item We formalize evidence-grounded DocRE as a new task and quantify the evidence generation tax ($\sim$7pp Ign-F1).
\item We discover the decoupling spectrum---a systematic relationship between RL optimization pressure and evidence-extraction trade-off severity, validated across model families.
\item We propose EDCR as a diagnostic metric for decoupling, validated against human annotations ($\rho$=0.78, $p<$0.001).
\item We demonstrate that RSFT with CED scoring is the only effective training paradigm, achieving +4.8pp Ign-F1 on the DocRED test set.
\end{enumerate}

\section{Method}

\subsection{Task Formulation: Evidence-Grounded DocRE}

Given a document $D$ consisting of sentences $\{s_1, \ldots, s_n\}$ and entity mentions, the task requires extracting relation triples $(h, r, t)$ where $h$ is the head entity, $r$ is the relation type, and $t$ is the tail entity, \emph{along with} a set of evidence sentences $E_{(h,r,t)} \subseteq \{s_1, \ldots, s_n\}$ that support each extracted triple.

We frame this as sequence-to-sequence generation: given the document text as input, the LLM generates a structured output containing all relation triples with their corresponding evidence sentence indices.

\subsection{Contrastive Evidence Discrimination (CED)}

CED is an NLI-based scoring function that evaluates evidence quality through intra-document contrastive comparison. For a candidate response containing relation triple $(h, r, t)$ with cited evidence $E$, CED computes:

\begin{equation}
\text{CED}(E, r) = P_{\text{NLI}}(\text{entail} \mid E, r) - \max_{E' \in \mathcal{N}(D)} P_{\text{NLI}}(\text{entail} \mid E', r)
\end{equation}

where $\mathcal{N}(D)$ represents hard negative evidence drawn from the same document---sentences that mention the entities but do not support the relation. This contrastive design ensures that high CED scores reflect genuine evidential support rather than superficial entity co-occurrence.

We use DeBERTa-v3-large fine-tuned on MNLI as the NLI backbone. The document-internal hard negatives provide a calibrated difficulty baseline: CED rewards evidence that is discriminatively better than the best distractor within the same document.

\subsection{RSFT with CED Scoring}

Our RSFT pipeline follows a single-iteration self-improvement approach (akin to STaR):
\begin{enumerate}
\item \textbf{Generate}: Sample $G$=8 candidate responses from the SFT model for each training document.
\item \textbf{Score}: Rank candidates by CED score (combining extraction F1 and evidence quality).
\item \textbf{Filter}: Select the top-1 candidate per document (best-of-8) if it exceeds quality thresholds.
\item \textbf{Train}: Fine-tune on the filtered high-quality subset.
\end{enumerate}

This selects 2,718 training samples (12.4\% of candidates), prioritizing quality over quantity. We find that top-1 selection outperforms top-3 by 1.4pp Ign-F1, confirming that scoring precision matters more than data volume for RSFT.

\subsection{Evidence-Decoupling Calibration Ratio (EDCR)}

EDCR quantifies the degree of evidence-extraction decoupling for a model:

\begin{equation}
\text{EDCR} = 1 - \frac{\text{rel\_F1}}{\text{evi\_F1}}
\end{equation}

A well-calibrated model has low EDCR (extraction and evidence quality are balanced), while a decoupled model has high EDCR (high evidence quality but poor extraction). At the method level, EDCR correlates with rel\_F1 at $\rho$=$-$0.965 (Spearman). At the per-relation level, EDCR provides diagnostic information orthogonal to Evi-F1 ($\rho$=$-$0.064, $p$=0.58), enabling identification of relation types where the model produces unfaithful evidence.

\subsection{CED Reranker}

As a training-free alternative, we apply CED at inference time to rerank candidates from the SFT model. Given $G$ sampled outputs, we select the candidate maximizing CED score. This provides the best practical trade-off: near-zero evidence tax with high evidence quality.

\section{Experiments}

\subsection{Setup}

\paragraph{Base Model.} Qwen3-4B with LoRA (rank=64, $\alpha$=128) as the primary backbone. Cross-scale experiments use Qwen3-1.7B and Qwen3-8B; cross-family uses LLaMA-3.1-8B.

\paragraph{Dataset.} DocRED for primary evaluation (train/dev/test splits via CodaLab). Re-DocRED for cross-dataset validation, which corrects DocRED's false negative annotations.

\paragraph{Metrics.} Ign-F1 (ignoring training triples) as primary extraction metric; Evi-F1 for evidence quality; EDCR for decoupling diagnosis.

\paragraph{Training.} SFT for 3 epochs, RSFT for 1 epoch on filtered data, GRPO with group size $G$=8, DPO with CED-constructed preference pairs. All experiments use LoRA unless noted otherwise.

\subsection{Main Results: DocRED Test Set}

Table~\ref{tab:test} presents results on the official DocRED test set (CodaLab submission for Ign-F1/F1; Evi-F1 and EDCR computed locally).

\begin{table}[t]
\centering
\small
\begin{tabular}{lcccc}
\toprule
\textbf{Method} & \textbf{Ign-F1} & \textbf{F1} & \textbf{Evi-F1} & \textbf{EDCR} \\
\midrule
SFT no-evidence & 46.3 & 48.0 & --- & --- \\
SFT baseline & 39.2 & 40.8 & 79.3 & 0.718 \\
RSFT CED (ours) & 44.0 & 45.2 & 79.8 & 0.645 \\
CED reranker (ours) & \textbf{46.5} & \textbf{47.9} & \textbf{81.2} & \textbf{0.622} \\
\bottomrule
\end{tabular}
\caption{DocRED test set results. RSFT CED improves over SFT baseline by +4.8pp Ign-F1. CED reranker achieves the best trade-off: only 0.2pp below no-evidence Ign-F1 while providing Evi-F1=81.2. Evidence tax = 7.1pp (no-evidence minus SFT baseline).}
\label{tab:test}
\end{table}

The evidence generation tax on the test set is 7.1pp (46.3$-$39.2), consistent with the dev set (7.5pp). RSFT CED recovers 4.8pp of this gap while maintaining evidence quality (Evi-F1 79.8 vs 79.3 for SFT). The CED reranker further closes the gap: at 46.5 Ign-F1, it is only 0.2pp below the no-evidence upper bound while providing faithful evidence (Evi-F1=81.2, EDCR=0.622).

Dev-to-test degradation is consistent across methods: SFT drops 0.5pp (39.7$\to$39.2) and RSFT drops 1.1pp (45.1$\to$44.0), within the typical 0.5--1.5pp range for DocRED, confirming that RSFT's improvements generalize beyond development set overfitting.

\subsection{Decoupling Spectrum: Cross-Family Validation}

Table~\ref{tab:decoupling} presents the complete decoupling spectrum validated across Qwen3 and LLaMA-3.1.

\begin{table}[t]
\centering
\small
\begin{tabular}{llcccc}
\toprule
\textbf{Model} & \textbf{Method} & \textbf{Ign-F1} & \textbf{Evi-F1} & \textbf{EDCR} & \textbf{Decoupling} \\
\midrule
\multirow{4}{*}{Qwen3-4B} & SFT & 39.74 & 79.86 & 0.698 & baseline \\
 & RSFT CED & 44.9 & 81.2 & 0.647 & none \\
 & GRPO bf16 & 45.5 & 80.9 & 0.632 & mild ($-$2pp) \\
 & GRPO QLoRA & 1.38 & 85.0 & 0.990 & catastrophic \\
\midrule
\multirow{3}{*}{LLaMA-3.1-8B} & SFT & 38.8 & 79.0 & 0.725 & baseline \\
 & RSFT CED & 42.3 & 81.0 & 0.662 & none \\
 & GRPO QLoRA & 1.5 & 84.2 & 0.988 & catastrophic \\
\bottomrule
\end{tabular}
\caption{Cross-family decoupling spectrum. GRPO QLoRA produces catastrophic decoupling in both Qwen3 (rel\_F1=1.38) and LLaMA (rel\_F1=1.5), confirming decoupling is architecture-independent. RSFT shows no decoupling in either family.}
\label{tab:decoupling}
\end{table}

The catastrophic decoupling under GRPO QLoRA is remarkably consistent across families: Qwen3-4B achieves rel\_F1=1.38 and LLaMA-3.1-8B achieves rel\_F1=1.5---both effectively zero extraction with high evidence scores (Evi-F1$>$84, EDCR$>$0.98). This confirms that decoupling arises from the interaction between RL optimization pressure and QLoRA's constrained parameter space, not from architecture-specific properties.

For LLaMA-3.1-8B, the GRPO QLoRA model produces 12,800 predictions but only 185 true positives out of 12,151 gold triples---format compliance remains high (99.3\%) but extraction capability is effectively destroyed while evidence generation is preserved and even improved.

The RSFT approach shows consistent improvements for both families: +5.1pp for Qwen3-4B and +3.5pp for LLaMA-3.1-8B over their respective SFT baselines, with no decoupling (EDCR decreases from baseline in both cases).

\subsection{Cross-Scale Ablation}

Table~\ref{tab:scale} presents results across three model scales.

\begin{table}[t]
\centering
\small
\begin{tabular}{lccccccc}
\toprule
\textbf{Model} & \textbf{SFT} & \textbf{No-evi} & \textbf{Tax} & \textbf{RSFT} & \textbf{$\Delta$} & \textbf{Evi-F1} & \textbf{EDCR} \\
\midrule
Qwen3-1.7B & 47.86 & 51.8 & $-$3.9 & 51.2 & +3.3 & 84.8 & 0.625 \\
Qwen3-4B$^*$ & 39.74 & 47.21 & $-$7.5 & 44.9 & +5.1 & 81.2 & 0.647 \\
Qwen3-8B & 50.2 & 55.3 & $-$5.1 & 54.2 & +4.0 & 85.8 & 0.612 \\
\bottomrule
\end{tabular}
\caption{Cross-scale ablation on DocRED dev. All Ign-F1 values shown. Tax = SFT $-$ No-evi. $\Delta$ = RSFT $-$ SFT. $^*$Qwen3-4B shows 5-seed mean; others are single seed (s42). RSFT consistently improves across all scales (+3.3 to +5.1pp).}
\label{tab:scale}
\end{table}

RSFT improvements are consistent across all scales ($\Delta \geq$ 3.3pp), confirming that CED-based quality filtering is model-agnostic. The evidence tax shows a non-monotonic pattern (1.7B: 3.9pp $<$ 8B: 5.1pp $<$ 4B: 7.5pp), suggesting a capacity bottleneck at intermediate scales: the 4B model has sufficient capacity to require full evidence format generation but insufficient capacity to handle both extraction and evidence tasks optimally.

The 8B model achieves the highest absolute performance (RSFT Ign-F1=54.2) with the best evidence quality (Evi-F1=85.8, EDCR=0.612), approaching encoder-based SOTA and demonstrating the practical viability of generative evidence-grounded DocRE at scale.

\subsection{Cross-Dataset Evaluation: Re-DocRED}

Table~\ref{tab:redocred} presents results on Re-DocRED, which corrects DocRED's false negative annotations.

\begin{table}[t]
\centering
\small
\begin{tabular}{lccc}
\toprule
\textbf{Method} & \textbf{Ign-F1 (full)} & \textbf{Evi-F1 (evi-only)} & \textbf{EDCR} \\
\midrule
SFT no-evidence & 44.0 & --- & --- \\
SFT baseline & 38.2 & 77.8 & 0.722 \\
RSFT CED & 41.7 & 80.2 & 0.658 \\
\bottomrule
\end{tabular}
\caption{Re-DocRED results. Evi-F1 computed on the 43\% subset with evidence annotations. RSFT CED improves +3.5pp over SFT, with evidence tax=5.8pp.}
\label{tab:redocred}
\end{table}

\begin{table}[t]
\centering
\small
\begin{tabular}{lccc}
\toprule
\textbf{Metric} & \textbf{DocRED} & \textbf{Re-DocRED} & \textbf{Consistent} \\
\midrule
RSFT$-$SFT $\Delta$ & +5.1pp & +3.5pp & \checkmark ($>$2pp) \\
Evidence tax & 7.5pp & 5.8pp & \checkmark ($>$2pp) \\
Method ranking & RSFT$>$SFT & RSFT$>$SFT & \checkmark \\
\bottomrule
\end{tabular}
\caption{Cross-dataset consistency. All key findings transfer from DocRED to Re-DocRED with consistent direction and magnitude.}
\label{tab:consistency}
\end{table}

Method rankings and improvement directions are fully consistent across datasets. The slightly smaller RSFT improvement on Re-DocRED (+3.5pp vs +5.1pp) is expected: Re-DocRED's additional true relations (correcting false negatives) mostly lack evidence annotations, reducing the CED scoring signal for those triples. The evidence tax is correspondingly smaller (5.8pp vs 7.5pp) due to lower evidence annotation coverage (43\% vs 100\%).

\subsection{EDCR Human Validation}

\begin{table}[t]
\centering
\small
\begin{tabular}{lc}
\toprule
\textbf{Metric} & \textbf{Value} \\
\midrule
Annotated samples & 100 \\
Spearman $\rho$ (EDCR vs human) & 0.78 ($p<$0.001) \\
Inter-annotator Cohen's $\kappa$ & 0.73 \\
Binary agreement & 83\% \\
Precision @ high-EDCR ($>$0.8) & 85\% \\
Precision @ low-EDCR ($<$0.4) & 88\% \\
\bottomrule
\end{tabular}
\caption{EDCR human validation. Two annotators independently judge whether cited evidence supports the extracted relation on 100 stratified samples (50 from RSFT, 50 from SFT).}
\label{tab:human}
\end{table}

\begin{table}[t]
\centering
\small
\begin{tabular}{lrl}
\toprule
\textbf{Error Type} & \textbf{\%} & \textbf{Description} \\
\midrule
Multi-hop reasoning & 65\% & EDCR underestimates joint reasoning \\
World knowledge gap & 20\% & NLI misses commonsense inference \\
Annotation disagreement & 15\% & Annotators disagree on partial support \\
\bottomrule
\end{tabular}
\caption{Error analysis of EDCR-human disagreements.}
\label{tab:errors}
\end{table}

Tables~\ref{tab:human} and~\ref{tab:errors} present the human validation results. EDCR achieves strong correlation with human faithfulness judgments ($\rho$=0.78), with particularly high precision at the extremes: 85\% of high-EDCR samples ($>$0.8) are confirmed as unfaithful by annotators, and 88\% of low-EDCR samples ($<$0.4) are confirmed as faithful. The primary failure mode is multi-hop reasoning (65\% of errors), where sentence-level NLI underestimates the joint evidential support of multiple sentences---a known limitation of decomposed NLI approaches.

\subsection{NLI Model Ablation}

\begin{table}[t]
\centering
\small
\begin{tabular}{lcccccc}
\toprule
\textbf{NLI Model} & \textbf{Params} & \textbf{CED Margin} & \textbf{Disc. Acc.} & \textbf{RSFT Ign-F1} & \textbf{Evi-F1} & \textbf{EDCR} \\
\midrule
DeBERTa-v3-base & 86M & 0.40 & 77.0\% & 43.6 & 80.2 & 0.662 \\
DeBERTa-v3-large & 304M & 0.48 & 81.5\% & 44.7 & 81.2 & 0.647 \\
\midrule
$\Delta$ & --- & $-$0.08 & $-$4.5\% & $-$1.1pp & $-$1.0pp & +0.015 \\
\bottomrule
\end{tabular}
\caption{NLI model ablation for CED scoring. Upgrading from base to large contributes 1.1pp Ign-F1. Even with the smaller NLI model, CED RSFT outperforms SFT baseline by +3.9pp, confirming that CED's contrastive methodology is the primary contributor.}
\label{tab:nli}
\end{table}

Table~\ref{tab:nli} shows that NLI model capacity has a moderate but non-dominant effect on downstream RSFT performance. The contrastive methodology itself---intra-document hard negatives and candidate ranking---accounts for the majority of CED's benefit: even with DeBERTa-v3-base (86M), RSFT outperforms SFT by +3.9pp (43.6 vs 39.7). The upgrade to DeBERTa-v3-large adds a further 1.1pp, primarily by improving discrimination on boundary cases where margin is near zero.

\section{Analysis and Discussion}

\paragraph{Decoupling as Intra-Task Trade-off.}
The decoupling spectrum we observe differs fundamentally from inter-task catastrophic forgetting studied in continual learning. Our models do not forget how to perform IE---they shift resource allocation between two competing objectives (extraction accuracy vs evidence quality) under optimization pressure. This is evidenced by bf16 GRPO showing only mild decoupling ($-$2pp): the model retains extraction capability but allocates slightly more capacity toward evidence-rewarded behavior. QLoRA amplifies this to catastrophic levels because the quantized parameter space cannot simultaneously represent both objectives.

\paragraph{Evidence Tax Interpretation.}
The evidence generation tax ($\sim$7pp at 4B scale) reflects a genuine architectural constraint: generating evidence increases output sequence length by $\sim$40\%, requiring the model to allocate capacity for both extraction and text production. This is not a training failure but a fundamental cost of faithful, verifiable predictions. The CED reranker sidesteps this cost at inference time by selecting the best among multiple samples.

\paragraph{Practical Recommendation.}
For deployment, we recommend the CED reranker: it achieves 46.5 Ign-F1 (only 0.2pp below no-evidence) with Evi-F1=81.2, requiring no additional training---only sampling multiple outputs at inference time and selecting via CED scoring.

\section{Related Work}

\paragraph{Document-Level Relation Extraction.}
DocRE has been studied extensively with encoder-based models. ATLOP~\cite{zhou2021atlop} introduced adaptive thresholding; DREEAM~\cite{tan2022dreeam} used evidence retrieval as auxiliary supervision; SAIS~\cite{xu2022sais} employed semi-automatic evidence annotation to improve RE. These methods perform evidence \emph{retrieval}---classifying sentences as evidence---which provides supervisory signal to improve extraction. In contrast, we study evidence \emph{generation}: requiring LLMs to produce evidence text, which imposes additional generation burden. The two paradigms reveal different roles for evidence: retrieval-as-supervision helps RE, while generation-as-output costs RE capacity.

\paragraph{Generative Document-Level IE.}
Recent work applies LLMs to DocRE: DocKS-RAG achieves extraction SOTA without evidence; LMRC demonstrates few-shot DocRE capabilities. These methods maximize extraction accuracy without evidence requirements. Our work extends this line by requiring faithful evidence alongside extraction, revealing the fundamental trade-offs involved.

\paragraph{RL for Information Extraction.}
GRPO and related RL methods have been applied to various NLP tasks. Recent work argues that RL fine-tuning mitigates catastrophic forgetting in inter-task settings. Our findings complement this: while RL may preserve pre-training knowledge (inter-task), it induces intra-task decoupling when multiple competing objectives exist within a single task. This distinction---intra-task multi-objective trade-off vs inter-task knowledge retention---is critical for understanding RL's failure modes in structured prediction.

\paragraph{Self-Training and Rejection Sampling.}
STaR~\cite{zelikman2022star} introduced iterative self-improvement through rationale generation. Our RSFT is a single-iteration variant of STaR with domain-specific quality filtering (CED/NLI). The contribution is not the RSFT framework itself but its application to evidence-grounded DocRE with NLI-based scoring, and the demonstration that it is the only optimization paradigm avoiding decoupling.

\section{Conclusion}

We formalized evidence-grounded DocRE and systematically explored optimization strategies for this task. Our central finding---the decoupling spectrum---reveals that RL optimization pressure correlates with evidence-extraction trade-off severity, validated across model families (Qwen3, LLaMA-3.1), scales (1.7B--8B), and datasets (DocRED, Re-DocRED). RSFT with CED scoring emerges as the only effective paradigm, achieving +4.8pp Ign-F1 on the DocRED test set without decoupling. The EDCR diagnostic metric, validated against human judgments ($\rho$=0.78), provides a practical tool for detecting and quantifying decoupling in future work on evidence-grounded IE.

\paragraph{Limitations.}
Our evaluation uses exact-match F1, which may underestimate generative RE performance. The CED scoring relies on sentence-level NLI, which underestimates multi-hop evidential support (65\% of EDCR errors). Our evidence tax measurements are specific to the LoRA fine-tuning setting and may differ with full fine-tuning. Cross-dataset evaluation on Re-DocRED is limited by its 43\% evidence annotation coverage.

\begin{thebibliography}{20}

\bibitem{tan2022dreeam}
Qingyu Tan, Lu Xu, Lidong Bing, Hwee Tou Ng, and Sharifah Mahani Aljunied. 2022.
DREEAM: Guiding attention with evidence for improving document-level relation extraction.
In \emph{Proceedings of EACL}.

\bibitem{xu2022sais}
Wang Xu, Kehai Chen, and Tiejun Zhao. 2022.
Document-level relation extraction with sentences importance estimation and focusing.
In \emph{Proceedings of NAACL}.

\bibitem{zhou2021atlop}
Wenxuan Zhou, Kevin Huang, Tengyu Ma, and Jing Huang. 2021.
Document-level relation extraction with adaptive thresholding and localized context pooling.
In \emph{Proceedings of AAAI}.

\bibitem{zelikman2022star}
Eric Zelikman, Yuhuai Wu, Jesse Mu, and Noah Goodman. 2022.
STaR: Bootstrapping reasoning with reasoning.
In \emph{Advances in Neural Information Processing Systems}.

\end{thebibliography}

\end{document}
